\section{Introduction}

Process verifiers are increasingly used to rank, filter, and allocate compute
over reasoning traces. Yet the standard reading of verifier scores is often too
optimistic. In answer-visible settings, a high score can reflect at least two
very different behaviors: sensitivity to local reasoning validity, or
sensitivity to whether the trace is simply consistent with a correct-looking
final answer. Ordinary discrimination metrics do not separate these behaviors,
so strong average performance can coexist with brittle counterfactual behavior.

This problem matters because verifier-guided reasoning has become a standard
pattern in mathematical reasoning and process supervision
\citep{cobbe2021training,lightman2023verify,wang2024mips}. At the same time,
recent work on reward models and LLM evaluators shows that strong average
judging performance can mask important evaluator biases
\citep{lambert2024rewardbench,panickssery2024selfprefer,wang2024selftaught}.
This work connects these threads by asking a narrower question: when a process
verifier can see the final answer, how much of its apparent process skill is
actually answer-sensitive shortcutting?

We treat answer-visible process verification as an auditing problem. Rather
than proposing a large new verifier architecture, we build an evaluation
protocol that makes shortcut reliance falsifiable. The protocol combines
executable matched-quartet benchmarks, answer-matched counterfactual
discrimination (\amcd), answer-swap sensitivity (\ass), and strong
same-backbone masked baselines. Minimal repair objectives serve as diagnostic
interventions: they test whether answer-sensitive behavior can be reduced
without simply deleting useful reasoning signals.

The empirical picture is narrower but more useful than a new-method story.
Ordinary metrics overstate the trustworthiness of answer-visible verification,
and \amcd{} is much more informative about downstream behavior than ordinary
AUROC. The strongest current deployment strategy is to remove direct answer
access. A secondary cleaner audit benchmark reaches the same qualitative
conclusion in a different regime: masking still improves process-facing
metrics even when utility no longer separates. A canonical external-source
\textsc{ProcessBench} package corroborates the same conservative read on real
step-level traces. There, the strongest current verifier is a simple step-only
frozen baseline, answer-only swaps still move visible scores much more than
masked scores, and even a fully constructed local-repair subset does not
produce a visible advantage. Taken together, these results support an
audit-and-deployment story, not a universal repair story.
Table~\ref{tab:claim_boundary} states the intended claim boundary explicitly.

\paragraph{Contributions.}
Our contributions are fourfold:
\begin{enumerate}[leftmargin=1.5em]
\item We introduce an audit protocol for answer-visible process verification
built around \amcd, \ass, and strong same-backbone masked baselines rather than
ordinary discrimination alone.
\item We show that the strongest current deployment configuration is a masked
reranker, while judge-style evaluators and visible rerankers remain
substantially answer-sensitive.
\item We provide mechanism evidence that local answer sensitivity tracks worse
utility in a pattern consistent with lower local process discrimination,
rather than merely co-occurring with weak utility.
\item We develop a benchmark package with two complementary regimes:
\textsc{Craft-Deploy}, a deployment-oriented naturalized generated benchmark,
and \textsc{Craft-Clean}, a cleaner audit benchmark accepted under an
artifact-focused acceptance criterion; and we corroborate the main qualitative
conclusions on a canonical external-source \textsc{ProcessBench} package with
trace, prefix, answer-swap, and local-repair views.
\end{enumerate}

\begin{table}[t]
\centering
\caption{Claim boundary for the paper. The manuscript is organized around
auditing, benchmark design, and deployment guidance rather than a new dominant
verifier architecture.}
\label{tab:claim_boundary}
\small
\begin{tabularx}{\linewidth}{p{0.24\linewidth}p{0.34\linewidth}X}
\toprule
Claim & Direct evidence & Not claiming \\
\midrule
Ordinary verifier metrics can mislead in answer-visible settings. &
\amcd{} and \ass{} separate ordinary valid-vs.-invalid discrimination from
answer-sensitive shortcutting across judges, frozen heads, and rerankers. &
Not claiming that every answer-visible verifier is useless, or that ordinary
AUROC has no value at all. \\
Masked scoring is the strongest current deployment choice. &
On \textsc{Craft-Deploy}, masked reranking beats visible reranking on
\amcd{}, exploitability, robustness, and seed-averaged reproduction. &
Not claiming that masking is the only possible long-run solution, only that it
is the safest current deployment choice in our setting. \\
Minimal repairs are diagnostic interventions rather than the final answer. &
\texttt{visible\_cond\_swap} moves the faithfulness--utility frontier, but no
simple visible repair dominates strong masked baselines or masked reranking. &
Not claiming universal repair dominance or a new best verifier architecture. \\
\textsc{Craft-Clean} is a usable cleaner audit benchmark. &
The accepted export keeps \texttt{52/54} families under an artifact-focused
criterion and reproduces the visible-vs.-masked faithfulness gap at mid-scale. &
Not claiming full benchmark closure; external human blind audit is still
pending. \\
External-source corroboration supports the benchmark pivot. &
On \textsc{ProcessBench}, the current strongest verifier is a step-only frozen
baseline, answer-only swaps still expose a large visible-vs.-masked gap, and a
fully constructed local-repair subset still does not favor the visible view. &
Not claiming that external-source results already support a new deployment
headline, or that \textsc{ProcessBench} is itself an \amcd/\ass{} quartet benchmark. \\
\bottomrule
\end{tabularx}
\end{table}
